\documentclass[11pt]{article}
\usepackage[margin=1in]{geometry}
\usepackage{amsmath,amssymb,amsthm}
\usepackage[round]{natbib}
\usepackage{hyperref}
\usepackage{microtype}

\newtheorem{theorem}{Theorem}
\newtheorem{lemma}{Lemma}
\newtheorem{remark}{Remark}

\title{A Feynman--Wigner-Style Diagnostic for the Efficacy\\
of Conformal Prediction via Signed de Finetti Representations}
\author{Peter Cotton}
\date{Draft --- companion note}

\begin{document}
\maketitle

\begin{abstract}
Conformal prediction builds prediction sets that cover the truth at a rate you choose,
finite-sample and distribution-free, assuming only exchangeable data. That
guarantee is marginal. de Finetti's theorem describes the exchangeability it rests on, and in
the finite form the mixing measure may be signed \citep{kerns2006}. We present a short lemma that decomposes the
slope of conformal's calibration-conditional coverage into two terms. One is a non-negative
threshold term, the classical Beta-law fan. The other carries the sign of the de Finetti
measure. Positive (extendable) mixtures make conformal conditionally adaptive. The signed
corner---de-meaned, ranked, or compositional scores---makes it anti-adaptive. The marginal
guarantee is the same either way. Only what it hides changes.
\end{abstract}

\section{Setup}

Split conformal returns a set $C(X)$ with $\mathbb P\big(Y\in C(X)\big)\ge 1-\alpha$ for any
distribution and any sample size, provided the data are exchangeable. That marginal guarantee
is the whole appeal, and it is not in dispute here. The question is what it says about the case
in front of you.

Scores $S_1,\dots,S_n$ (calibration) and $S_{n+1}$ (test) are jointly
exchangeable with a continuous joint law (no ties). The conformal width is the order
statistic $\hat q = S_{(k)}$ of the calibration scores, $k=\lceil(n+1)(1-\alpha)\rceil$, and a
test point is covered when $S_{n+1}\le\hat q$.

Two kinds of dependence live in any such problem, and they should not be confused.
\begin{itemize}
\item Within a row: inside a single $(x,y)$, does the residual's shape depend on $x$? That is
the subject of the companion paper \citep{cotton_marginally}, where the log-score regret of a
single-shape conformal system equals the mutual information $I(R;X)$. It is a property of one
observation's law, present even when the rows are i.i.d.
\item Across rows: how do distinct observations relate to one another, whether independent,
sharing a latent cause, or mutually constrained? That is the subject of de Finetti's theorem,
and of this note.
\end{itemize}
The two axes are orthogonal. This note is entirely about the second.

\section{Finite de Finetti and the sign of the mixture}

de Finetti's theorem pins down what infinite exchangeability is.

\begin{theorem}[de Finetti; \citealp{definetti1937,hewitt1955}]\label{thm:definetti}
Let $X_1,X_2,\dots$ be an infinitely exchangeable sequence in a Polish space. Then there is a
probability measure $\pi$ on the space of distributions such that, for every $n$, the law of
$(X_1,\dots,X_n)$ equals $\int P^{\otimes n}\,\mathrm{d}\pi(P)$. Equivalently, conditional on a
random measure $\Theta\sim\pi$, the $X_i$ are i.i.d.\ with law $\Theta$.
\end{theorem}

The mixing measure $\pi$ is a genuine prior, and that forces a sign. A positive mixture has
$\operatorname{Cov}(S_i,S_j)=\operatorname{Var}\big(\mathbb E[S_i\mid\Theta]\big)\ge 0$, so de
Finetti mixtures are non-negatively correlated. Finite exchangeable sequences are not so
constrained. Exchangeability alone forces only $\rho\ge -1/(n-1)$ \citep{aldous1985}, and the
negative band is real---sampling without replacement sits at its floor.

The finite theory completes the picture. \citet{diaconis1980} show a finite exchangeable
sequence is approximately a positive i.i.d.\ mixture, within total variation $O(k/n)$.
\citet{kerns2006} make it exact, at the price of signing.

\begin{theorem}[Kerns--Sz\'ekely; \citealp{kerns2006}]\label{thm:signed}
Let $(X_1,\dots,X_n)$ be exchangeable on a measurable space. Then there is a finite
signed measure $\mu$ of total mass $1$ on the space of distributions with
$\mathrm{law}(X_1,\dots,X_n)=\int P^{\otimes n}\,\mathrm{d}\mu(P)$. One may take $\mu\ge 0$ (a
genuine prior, recovering Theorem~\ref{thm:definetti}) if and only if the sequence extends to an
infinitely exchangeable one \citep{leonetti2018}.
\end{theorem} We picture this as the i.i.d.\ laws tracing a curve: positive mixtures are its convex hull, the
extendable $\rho\ge 0$ laws, and signed mixtures its affine span, everything. A signed
representation marks non-extendability---a finite exchangeable sequence that cannot be continued
to an infinite one---which is the negative-correlation regime.

We frame the signed measure as a negative probability, the kind \citet{feynman1987} used
partway through a calculation, or \citet{wigner1932}'s phase-space distribution, which dips
below zero while every observable marginal stays a genuine probability (the idea reached
statistics through \citealp{bartlett1945}). You cannot sample a parameter from it. That is why
the signed corner gives a diagnosis, not a model: recovering a latent means fitting a real one,
not reading the quasi-probability off.

So the sign of the de Finetti mixing measure is a clean label: positive = extendable,
non-negatively correlated, a genuine prior. Signed = finite, negatively correlated,
no genuine prior.

\section{Marginal coverage ignores the sign}

For continuous exchangeable scores the rank of $S_{n+1}$ among all $n+1$ is uniform, and
$S_{n+1}\le S_{(k)}$ iff that rank is at most $k$, so
\begin{equation}
\mathbb P\big(S_{n+1}\le\hat q\big)=\frac{k}{n+1}\ge 1-\alpha .
\end{equation}
This uses nothing but exchangeability \citep{vovk2005}. Positive or signed, the marginal
number is the same. The sign enters only when we ask what the certificate hides.

\section{The conditional-coverage decomposition}

Condition on the realised width. For a value $q$ of the width $\hat q$ and a threshold $s$, write
$H(s,q)=\mathbb P(S_{n+1}\le s\mid\hat q=q)$. Coverage holds when $S_{n+1}\le\hat q$, so coverage
conditional on the realised width is the diagonal $c(q)=H(q,q)$, where the threshold equals the
width.

\begin{lemma}[Conditional-coverage slope]\label{lem:slope}
If $(S_{n+1},\hat q)$ has a joint density, then
\begin{equation}
c'(q)=\underbrace{f_{S_{n+1}\mid\hat q}(q\mid q)}_{\textup{(i) threshold}}
\;+\;\underbrace{\partial_q\,\mathbb P(S_{n+1}\le q\mid\hat q=q)}_{\textup{(ii) dependence}} .
\end{equation}
Term \textup{(i)} is non-negative. Term \textup{(ii)} is $\le 0$ when $S_{n+1}$ is
stochastically increasing in $\hat q$, $=0$ under independence, and $\ge 0$ when stochastically
decreasing.
\end{lemma}

\begin{proof}
Total derivative of $H(q,q)$ along the diagonal: $c'(q)=\partial_s H(s,q)\big|_{s=q}+\partial_q
H(s,q)\big|_{s=q}$. The first equals $f_{S_{n+1}\mid\hat q}(q\mid q)\ge 0$. For the second,
$\partial_q H(s,q)\le 0$ for all $s$ exactly when $\mathbb P(S_{n+1}\le s\mid\hat q=q)$ is
non-increasing in $q$, the definition of $S_{n+1}$ being stochastically increasing in
$\hat q$. The reverse inequality is stochastic decrease.
\end{proof}

\begin{remark}[The independence case is classical]
Under independence term (ii) vanishes and $c(q)=F(q)$, the test-score CDF, which already rises
with the width purely from estimation. This is the known law of calibration-conditional
coverage: $F(\hat q)=F(S_{(k)})\sim\operatorname{Beta}(k,n-k+1)$
\citep{vovk2012,marques2024,angelopoulos2023gentle}. Lemma~\ref{lem:slope} is its extension to
dependent scores, with term (ii) the new piece.
\end{remark}

\begin{remark}[Reading the sign]
The sign of the cross-sample dependence sets the direction. Positive dependence makes conformal
conditionally adaptive, negative dependence makes it anti-adaptive. The mechanism is short. The
width $\hat q=S_{(k)}$ is non-decreasing in the calibration scores, and the test score is
disjoint from them. A positive (proper) de Finetti mixture makes the score vector associated
\citep{esary1967}, so $\hat q$ and $S_{n+1}$ move together: term (ii) is $\le 0$, it cancels the
threshold term, and pooling the calibration scores conditions you into the realised world. A
signed, negatively-correlated law makes the vector negatively associated \citep{joagdev1983}
(without-replacement samples, multinomials, ranks, compositional vectors), so term (ii) is
$\ge 0$ and adds to it: the calibration set anti-predicts the test, and the per-case coverage
swings high and low while the marginal stays pinned. The order-statistic version uses a
difference in place of the derivative and behaves the same.
\end{remark}

The reading is a go/no-go on conformal's coverage, and coverage only. The residual-information
gap of the companion paper sits there whatever the sign, so even at its best this is an honest
certificate, not a sharp forecast. With independent or positively-dependent data the certificate
is at least conditionally honest. With negatively-dependent data it is not: the number holds on
average while the per-case coverage swings high and low. The sharpness comes from modelling the
latent, never from the wrapper.

\section{Why the signed corner is not exotic: constrained scores}

The negative corner arises whenever the nonconformity score is defined relative to the
batch, because a linear constraint forces it. De-meaning ($\sum_i s_i=0$), differencing,
ranks or percentiles-within-batch, and compositional or budgeted targets all impose such a
constraint, and the constraint is the negative association \citep{joagdev1983}. For
these scores conformal remains marginally valid---if the relativisation is symmetric across
calibration and test, exchangeability is preserved---but its conditional coverage is the
fanning, anti-adaptive kind of Lemma~\ref{lem:slope}.

A caution: if the relativisation treats the test point asymmetrically (de-meaning by a
training-only mean, say), exchangeability itself fails and even the marginal guarantee is no
longer assured. This is the familiar reason split conformal uses out-of-fold rather than
in-sample residuals \citep{lei2018}; in-sample regression residuals are both constrained
($\sum_i e_i=0$) and heteroscedastic (variance $1-h_{ii}$), and so are not exchangeable to
begin with.

\begin{remark}[Rank, choice, and market-implied scores]
The signed corner is the native habitat of contest data. In a Thurstone or Luce choice model a
field of competitors carries latent values and the observed outcome is a winner or a full
ranking. The winner indicator is a single multinomial trial, and the rank vector a
permutation law, both negatively associated \citep{joagdev1983}. Any nonconformity score
read off such a field---a rank, a within-field relative score, or a market-implied win
probability---therefore sits in the negatively-associated regime above, where the marginal
certificate is least informative about conditional reliability. The converse direction is the
modelling the certificate cannot supply: recovering latent means from observed prices or
winning probabilities \citep{cotton_horserace} is exactly the de Finetti latent-recovery of
Section~2. In betting terms the market price is the crowd's $Q$, the recovered latent mean is an
estimate of the true $P$, and the divergence between them is the rent of the companion paper.
So the same contest data that lands conformal in its anti-adaptive corner is also where the de
Finetti latent, once estimated, pays.
\end{remark}

\section{A runnable detector}

The gap is defined through the log-score, so it is not something you read off data. Its
energy-statistics counterpart is, and it gives a test you can run on any fitted conformal
predictor.

\begin{lemma}[A distance-covariance detector]\label{lem:dcov}
Let $U=\widehat G(R)$ be the conformal rank of the residual among the calibration residuals.
Then $U$ is marginally uniform, and $\operatorname{dCov}(U,X)=0$ iff $U\perp X$ iff $R\perp X$ iff
$I(R;X)=0$. The permutation test that reshuffles the $U$ labels against $X$ is a valid test of
$R\perp X$, and rejecting it certifies $I(R;X)>0$: the conformal rank is conditionally
non-uniform, so the marginal coverage is uneven across $x$.
\end{lemma}

Two reasons this is the right object. $U$ is standardized to uniform whatever the model, so the
test probes only the conditional structure conformal cannot see. And distance covariance
\citep{szekely2007} catches dependence in any moment: under symmetric heteroscedastic noise the
conditional mean of $U$ stays at $\tfrac12$ while its spread moves with $x$, so a correlation
test sees nothing and the distance covariance sees it.

A small experiment makes the point (one run, reproduced by \texttt{check\_dcov.py}). The
residual is $R=\sigma(X)\,\varepsilon$ with the location correct, so any dependence is pure
conditional shape.

\begin{center}
\begin{tabular}{lcccc}
\hline
regime & marginal coverage & $\operatorname{dCov}(U,X)$ & permutation $p$ & Pearson \\
\hline
homoscedastic & $0.90$ & $0.001$ & $0.33$ & $+0.02$ \\
heteroscedastic, mild & $0.89$ & $0.011$ & $0.002$ & $+0.05$ \\
heteroscedastic, strong & $0.90$ & $0.023$ & $0.002$ & $-0.03$ \\
\hline
\end{tabular}
\end{center}

Marginal coverage is on target throughout, so it cannot separate the regimes. The test fires
only where the gap is positive, and the Pearson column shows that a correlation test would miss
it.

The idea of a dependence measure as a conditional-coverage diagnostic is due to
\citet{feldman2021}, who use HSIC between the miscoverage indicator and the interval length, and
\citet{braun2025}, who classify the indicator on $x$. Testing the rank against the features is
the rank-based, level-free variant, and it is the energy counterpart of the gap: $I(R;X)$ is the
dependence of $R$ and $X$ in KL \citep{cover2006}, and the distance covariance, equivalently
HSIC \citep{sejdinovic2013}, is the same dependence in energy. The paper's gap is the one you
prove; this is the one you run.

It detects, it does not certify. A large $p$ means nothing was found at this power, not
conditional validity, by the no-go results discussed in the companion paper
\citep{cotton_marginally}. The ranks share the calibration ECDF, so leave-one-out or
full-conformal ranks make the null exact. The test ships as \texttt{conformalguide.guardrails}.

\section{Related work}

The marginal-coverage fact is standard \citep{vovk2005}. The Beta law for coverage given the
calibration set is \citet{vovk2012}, refined by \citet{marques2024}; training-conditional
coverage is \citet{bian2023}. Magnitude-based bounds for dependent data are
\citet{barber2023beyond} and \citet{oliveira2024}. Association and negative association are
\citet{esary1967} and \citet{joagdev1983}; finite and signed de Finetti are
\citet{diaconis1980} and \citet{kerns2006,leonetti2018}.

\citet{barberpananjady2026} are closest: they show conformal can over- or under-cover under
temporal dependence. That account is marginal, it grades a departure by the size of the
dependence rather than its sign, and it does not condition on the calibration set. The sign,
and its identification with the sign of the finite de Finetti measure, is what this note adds.

de Finetti describes a law. Conformal runs without one. Using a de Finetti measure means
estimating it, and that estimation is ill-posed: the latent is recovered from finite data, so in
general one regularizes, shrinking toward the field. Conformal skips this, and the
residual-information gap is the price of skipping it. A signed measure gives nothing to estimate
from at all, so the representation tells you how conformal behaves rather than standing in for
it. \citet{datta2025} argue conformal is not a de Finetti conditional at all.

\bibliographystyle{plainnat}
\bibliography{../notes}

\end{document}
